% Lösungen Einheit 3
\einheitenkopf[][3 Krümmung und Wendepunkte]{Einheit 3 von 5 · Krümmung und Wendepunkte}

\erg{23}{a) $f''(x) = 6x - 18$, $f'''(x) = 6$ \quad b) $f''(x) = -12x + 10$, $f'''(x) = -12$ \quad c) $f''(x) = -12x^2 + 6$, $f'''(x) = -24x$ \quad d) $f''(x) = 3x^2 - 6x$, $f'''(x) = 6x - 6$ \quad e) $f''(x) = 6x - 12 = 0$ für $x = 2$, $f'''(2) = 6 \ne 0$: $W(2\,|\,0)$ \quad f) $f''(x) = 12x^2 - 12x = 12x\,(x - 1)$; $f'''(x) = 24x - 12$, $f'''(0) = -12 \ne 0$, $f'''(1) = 12 \ne 0$: $W_1(0\,|\,1)$, $W_2(1\,|\,0)$ \quad g) $f'(x) = -x\,e^{-x}$, $f''(x) = (x - 1)\,e^{-x} = 0$ für $x = 1$; $f'''(x) = (2 - x)\,e^{-x}$, $f'''(1) = e^{-1} \ne 0$: $W(1\,|\,2e^{-1}) \approx W(1\,|\,0{,}74)$}

\erg{24}{a) $f''(x) = 6x + 6$: rechtsgekrümmt auf $]-\infty;\,-1]$, linksgekrümmt auf $[-1;\,\infty[$ \quad b) $f''(x) = -6x + 12$: linksgekrümmt auf $]-\infty;\,2]$, rechtsgekrümmt auf $[2;\,\infty[$ \quad c) $f''(x) = 12x^2 - 12$: linksgekrümmt auf $]-\infty;\,-1]$, rechtsgekrümmt auf $[-1;\,1]$, linksgekrümmt auf $[1;\,\infty[$}

\erg{25}{a) $f''(x) = 6x + 12$, $f''(-2) = 0$, $f'''(-2) = 6 \ne 0$ \quad b) $f''(x) = -6x$, $f''(0) = 0$, $f'''(0) = -6 \ne 0$ \quad c) $f''(x) = 6x^2 - 6$, $f''(1) = 0$, $f'''(1) = 12 \ne 0$, $f(1) = 0{,}5 - 3 + 2 = -0{,}5$ \quad d) $f''(x) = 6x - 6 = 0$ für $x = 1$, $f'''(1) = 6 \ne 0$, $W(1\,|\,2)$; $f'(1) = 1$: $w(x) = x + 1$ \quad e) $f'(x) = 4x^3 + 6x^2$, $f'(0) = 0$; $f''(x) = 12x^2 + 12x$, $f''(0) = 0$; $f'''(x) = 24x + 12$, $f'''(0) = 12 \ne 0$; $f(0) = 5$: Wendepunkt mit waagerechter Tangente, also Sattelpunkt \quad f) Links vom Tiefpunkt fällt der Graph aus der Nähe der $x$-Achse auf $-e^2$; weit links ist er dabei rechtsgekrümmt (er biegt sich von der Achse nach unten weg), am Tiefpunkt linksgekrümmt. Die Krümmung wechselt, also gibt es links von $x = 2$ einen Wendepunkt (Kontrolle: $f''(x) = (x - 1)\,e^{x}$, $W(1\,|\,{-}2e)$).}

\erg{26}{a) Emma vertauscht links und rechts: $f''(x) > 0$ heißt linksgekrümmt. Richtig: linksgekrümmt für $x \le 1$, rechtsgekrümmt für $x \ge 1$ \quad b) Paul nimmt $f''(-1) = 0$ als $y$-Wert. Richtig: $f(-1) = -1 + 3 = 2$, also $W(-1\,|\,2)$}

\erg{27}{a) Nein: $f''(x) = 12x^2$ ist links und rechts von 0 positiv, wechselt also das Vorzeichen nicht; der Graph ist überall linksgekrümmt, bei 0 liegt ein Tiefpunkt. \quad b) Im Wendepunkt wechselt die Krümmung: auf der einen Seite liegt der Graph über der Tangente (linksgekrümmt), auf der anderen darunter (rechtsgekrümmt).}

\erg{28}{a) $f''(x) = -0{,}6x + 1{,}2 = 0$ für $x = 2$, $f'''(x) = -0{,}6 \ne 0$, $f(2) = 1{,}6$: $W(2\,|\,1{,}6)$; für $x < 2$ ist $f'' > 0$ (Linkskurve), danach Rechtskurve – bei 200 m in $x$-Richtung und 160 m in $y$-Richtung \quad b) Im Wendepunkt ist die Krümmung null: die Kurve wechselt von links nach rechts, das Lenkrad geht dabei durch die Mittelstellung.}
